% Responsable: Aldair Herrera (Documentación) — 300 a 400 palabras
\section*{Resumen ejecutivo}
\addcontentsline{toc}{section}{Resumen ejecutivo}

En la Facultad de Ciencias de la Computación de la Universidad Técnica Estatal de
Quevedo, la selección de docentes no titulares se hace sin un sistema propio. Según las
entrevistas al Decano y a la Coordinadora de la carrera, los postulantes envían su hoja
de vida por correo a Talento Humano o directamente al decanato. Allí se descarga cada
una, se califica a mano con una matriz de valoración en Excel y, con los puntajes, se
redacta un informe que pasa por el Consejo Directivo, el Consejo Académico y el Consejo
Universitario. El Decano señaló que lo que más tiempo toma es revisar las hojas de vida
y registrar los puntajes, y la Coordinadora afirmó que el proceso tiene que
automatizarse.

El Sistema de Selección Docente y Convocatorias (SSDC) ya existe como un monolito en
capas desarrollado en la asignatura Aplicaciones Web. En este proyecto se evoluciona
hacia una arquitectura distribuida. Sus usuarios son los postulantes, que envían sus
documentos y consultan el estado de su postulación; el Decano y la Coordinadora, que
califican y generan el informe; y las autoridades que aprueban la contratación.

Un sistema distribuido es apropiado porque el SSDC depende de servicios externos, como
el correo, el almacenamiento de archivos y la inteligencia artificial, que hoy se
ejecutan dentro del mismo proceso que atiende postulaciones y evaluaciones. Si uno de
ellos falla, puede afectar todo el sistema, y cualquier cambio obliga a reiniciar el
backend completo. Al separar el sistema en servicios con su propia base de datos, cada
falla queda acotada al servicio que la produce y cada parte se despliega sin detener
las demás.

En esta primera entrega se definen once requisitos funcionales y seis no funcionales
relacionados con las entrevistas, la arquitectura base con el modelo C4 en sus niveles
de contexto y contenedores, el modelo de comunicación entre servicios, el análisis CAP y
PACELC de los datos críticos y los riesgos técnicos derivados de incidentes reales.
Además, se entrega un esqueleto con dos servicios, convocatorias y notificaciones, que
se comunican por HTTP y se levantan con un solo comando. En las entregas siguientes se
implementarán los microservicios y sus contenedores, la base de datos distribuida y el
sistema final con sus pruebas.